\documentclass[12pt]{article}
\usepackage[utf8]{inputenc}
\usepackage[T1]{fontenc}
\usepackage{hyperref}
\title{Paranan Language Resources for K-12 Curriculum Development}
\author{Compiled by Claude Sonnet 5.5 (claude-sonnet-5-5) for Andr\'as Kornai}
\date{}
\begin{document}
\maketitle

\section{Introduction}
Paranan (Palanan) is a Northern Luzon language of the Philippines spoken on the northeastern coast of Isabela, with about 15,000 speakers. It should not be confused with Pahanan (Paranan) Agta, a distinct Aeta language of the same area that shares much vocabulary.

\section{Primary Educational Resources}
\begin{itemize}
\item \emph{Ortograpiya ni Paranan} (2020), orthography guide by the Komisyon sa Wikang Filipino and DepEd Isabela, covering graphemes, syllabication, spelling and pronunciation: \url{https://kwfwikaatkultura.ph/wp-content/uploads/2024/03/Paranan.pdf}
\end{itemize}
No primer, textbook or reader was found.

\section{Linguistic and Reference Materials}
\begin{itemize}
\item Wikipedia, Paranan language: \url{https://en.wikipedia.org/wiki/Paranan_language}
\item Wikipedia, Paranan Agta language (related, distinct): \url{https://en.wikipedia.org/wiki/Paranan_Agta_language}
\end{itemize}
No dictionary or grammar was found.

\section{Religious and Community Resources}
None identified.

\section{Digital and Audio Resources}
None identified.

\section{Curriculum Development Context}
The language is widely spoken, and a standard orthography now exists, which is the main prerequisite for materials. Beyond that, the search found very little. Dictionary, reader and lesson development would start essentially from scratch, with Isabela DepEd and KWF as natural partners.

\end{document}